\documentclass[tikz,border=10pt]{standalone}
\usepackage{amsmath}
\usepackage{tikz}
\usetikzlibrary{shapes.geometric, arrows.meta, positioning, fit, backgrounds, shadows, decorations.pathmorphing}

\begin{document}

\tikzset{
    process/.style={
        rectangle,
        rounded corners=3mm,
        draw=blue!70,
        fill=blue!15,
        thick,
        minimum height=1.1cm,
        minimum width=4cm,
        text width=3.7cm,
        align=center,
        font=\small\bfseries,
        drop shadow
    },
    component/.style={
        rectangle,
        rounded corners=2mm,
        draw=green!70,
        fill=green!10,
        thick,
        minimum height=0.9cm,
        minimum width=3.5cm,
        text width=3.2cm,
        align=center,
        font=\scriptsize
    },
    detail/.style={
        rectangle,
        rounded corners=1.5mm,
        draw=orange!70,
        fill=orange!8,
        minimum height=0.65cm,
        minimum width=3.2cm,
        text width=3cm,
        align=center,
        font=\tiny
    },
    arrow/.style={
        -Stealth,
        very thick,
        draw=gray!70
    },
    bigarrow/.style={
        -Stealth,
        line width=1.5mm,
        draw=blue!40
    },
    label/.style={
        font=\scriptsize\bfseries,
        text=blue!70!black
    },
    title/.style={
        font=\Large\bfseries,
        text=blue!80!black
    },
    hamiltonian/.style={
        rectangle,
        rounded corners=2mm,
        draw=purple!70,
        fill=purple!10,
        thick,
        minimum height=2.8cm,
        minimum width=3.8cm,
        text width=3.5cm,
        align=left,
        font=\tiny
    }
}

\begin{tikzpicture}[node distance=0.7cm and 1.2cm]

    % Title
    \node[title] (title) at (0,0) {Monte Carlo Simulation Flow for Nanogranular Magnetic Structures};
    \node[font=\footnotesize\itshape, text=gray!70, below=0.1cm of title] (subtitle) {(Agudelo Giraldo 2016)};

    % Step 1: Structural Generation
    \node[process, below=0.6cm of subtitle] (struct) {1. Structural\\Generation};
    \node[component, below=0.2cm of struct] (struct1) {Lognormal volume\\distribution};
    \node[detail, below=0.15cm of struct1] (struct2) {Periodic boundary\\conditions (x-y plane)};

    % Step 2: Grain Boundary Relaxation
    \node[process, below=0.5cm of struct2] (relax) {2. Grain Boundary\\Relaxation};
    \node[component, below=0.2cm of relax] (relax1) {Morse Potential};
    \node[detail, below=0.15cm of relax1] (relax2) {Monte Carlo Method\\(structural)};

    % Step 3: Hamiltonian Definition
    \node[process, below=0.5cm of relax2] (ham) {3. Hamiltonian\\Construction};

    % Hamiltonian box (detailed)
    \node[hamiltonian, below=0.2cm of ham] (hambox) {
        \textbf{4 Components:}\\[0.1cm]
        (i) \textbf{Exchange:} $J_{ij}$\\
        \phantom{(i)} Distance-dependent\\[0.05cm]
        (ii) \textbf{Anisotropy:}\\
        \phantom{(i)} Magnetocrystalline, surface,\\
        \phantom{(i)} grain boundary\\[0.05cm]
        (iii) \textbf{External field:} $H_{\text{ext}}$\\[0.05cm]
        (iv) \textbf{Dipolar interaction:}\\
        \phantom{(i)} Spin-spin + spin-cell
    };

    % Step 4: Monte Carlo Relaxation
    \node[process, below=0.5cm of hambox] (mc) {4. Magnetic State\\Relaxation};
    \node[component, below=0.2cm of mc] (mc1) {Metropolis Algorithm};
    \node[detail, below=0.15cm of mc1] (mc2) {Classical Heisenberg\\+ Cubic anisotropy};

    % Step 5: Parallelization
    \node[process, below=0.5cm of mc2] (para) {5. Parallelization};
    \node[component, below=0.2cm of para] (para1) {Cell decomposition};
    \node[detail, below=0.15cm of para1] (para2) {Shared memory access\\Cutoff radius: 5 replicas};

    % Step 6: Results
    \node[process, below=0.5cm of para2, fill=green!20, draw=green!70] (results) {6. Results \& Analysis};
    \node[component, below=0.2cm of results, fill=yellow!10] (res1) {M(T), Hysteresis loops};
    \node[detail, below=0.15cm of res1] (res2) {ZFC-FC, Blocking temp.\\Mono/Multi-domain};

    % Right side: Key findings
    \node[process, right=1.5cm of ham, minimum width=4.5cm, text width=4.2cm, fill=purple!10, draw=purple!70] (findings) {Key Physical Effects};

    \node[component, below=0.2cm of findings, minimum width=4.2cm, text width=3.9cm, fill=purple!5] (find1) {Monodomain tendency:\\$\uparrow$ thickness, $\uparrow$ grain vol.};

    \node[component, below=0.25cm of find1, minimum width=4.2cm, text width=3.9cm, fill=purple!5] (find2) {Perpendicular $\to$ in-plane\\transition (surface vs dipolar)};

    \node[component, below=0.25cm of find2, minimum width=4.2cm, text width=3.9cm, fill=purple!5] (find3) {Superparamagnetic effect\\(boundary anisotropy isolation)};

    % Arrows
    \draw[arrow] (struct2) -- (relax);
    \draw[arrow] (relax2) -- (ham);
    \draw[arrow] (hambox) -- (mc);
    \draw[arrow] (mc2) -- (para);
    \draw[arrow] (para2) -- (results);

    % Connection to findings
    \draw[arrow, dashed, purple!70] (hambox.east) -- (findings.west);
    \draw[arrow, dashed, purple!70] (mc.east) -- ++(0.5,0) |- (find2.west);
    \draw[arrow, dashed, purple!70] (results.east) -- ++(0.5,0) |- (find3.west);

    % Temperature note
    \node[detail, right=1.5cm of relax, minimum width=4.2cm, text width=4cm, draw=red!60, fill=red!5] (tempnote) {Temperature-dependent\\anisotropy included};
    \draw[arrow, dashed, red!60] (tempnote.west) -- (relax.east);

    % Disorder note
    \node[detail, right=1.5cm of struct, minimum width=4.2cm, text width=4cm, draw=orange!70, fill=orange!8] (disnote) {Disorder: anisotropy vector\\$\propto$ lattice deformations};
    \draw[arrow, dashed, orange!70] (disnote.west) -- (hambox.north east);

\end{tikzpicture}

\end{document}
